\documentclass[11pt,a4paper]{article}
\usepackage[T1]{fontenc}
\usepackage[utf8]{inputenc}
\usepackage{lmodern,microtype}
\usepackage[margin=25mm,headheight=15pt]{geometry}
\usepackage{amsmath,amssymb,amsthm,mathtools}
\usepackage{booktabs,array,longtable,enumitem}
\usepackage[dvipsnames]{xcolor}
\usepackage{fancyhdr,titlesec,xurl}
\usepackage[colorlinks=true,linkcolor=MidnightBlue,citecolor=MidnightBlue,urlcolor=MidnightBlue]{hyperref}
\usepackage{bookmark}
\definecolor{ink}{RGB}{25,46,65}
\titleformat{\section}{\large\bfseries\color{ink}}{\thesection}{0.7em}{}
\titleformat{\subsection}{\normalsize\bfseries\color{ink}}{\thesubsection}{0.7em}{}
\pagestyle{fancy}\fancyhf{}
\fancyhead[L]{\small Signed quadratic-feedback inversion}
\fancyhead[R]{\small Research draft, 29 September 2026}
\fancyfoot[C]{\thepage}
\renewcommand{\headrulewidth}{0.3pt}
\setlength{\parskip}{3pt}
\setlength{\emergencystretch}{3em}
\setlist{itemsep=4pt,topsep=5pt,leftmargin=1.7em}
\numberwithin{equation}{section}
\newtheorem{theorem}{Theorem}[section]
\newtheorem{lemma}[theorem]{Lemma}
\newtheorem{proposition}[theorem]{Proposition}
\newtheorem{corollary}[theorem]{Corollary}
\theoremstyle{definition}
\newtheorem{definition}[theorem]{Definition}
\newtheorem{example}[theorem]{Example}
\newtheorem{question}[theorem]{Research question}
\theoremstyle{remark}\newtheorem{remark}[theorem]{Remark}
% ed. (2026-09-29): unnumbered environment for ProveIt editorial notes, so that the article's own numbering is unchanged
\newtheorem*{ednote}{Editorial note (ProveIt, 2026-09-29)}
\newcommand{\R}{\mathbb R}\newcommand{\C}{\mathbb C}\newcommand{\N}{\mathbb N}
\newcommand{\ee}{\mathrm e}\newcommand{\dd}{\mathrm d}
\newcommand{\U}{\widehat U}\newcommand{\V}{\widehat V}
\newcommand{\B}{\mathcal B}\newcommand{\I}{\mathcal I}
\newcommand{\coef}[2]{[#1^{#2}]}
\DeclareMathOperator{\Res}{Res}
\hypersetup{pdftitle={Signed Quadratic-Feedback Inversion: A Finite-Core Resolution and Sharp Entire Borel Growth},pdfauthor={AI-assisted research draft prepared for Vladimir Reshetnikov},pdfsubject={Resolution of a repository-local inverse asymptotic conjecture, signed finite-core universality, and Borel growth}}
\begin{document}
\begin{center}
{\LARGE\bfseries\color{ink} Signed Quadratic-Feedback Inversion\par}
\vspace{0.5em}
{\Large A finite-core resolution and sharp entire Borel growth\par}
\vspace{1em}
Research article prepared for Vladimir Reshetnikov\\
29 September 2026
\end{center}
\begin{abstract}
The ProveIt article on finite-core universality leaves conjectural a signed
coefficient equivalent for the compositional inverse of
$\U(q)=\sum_{j\ge1}q^j\exp(a j^2\U(q))$, $a>0$.
We give a proof by reversing the defining equation before performing
large-order analysis. The crucial estimate is an arbitrary-algebraic-accuracy
reduction to one action above a sufficiently large analytic core. Its proof
uses absolute bounds only after the cancellations within the core have been
resummed. A merger defect for quadratic actions suppresses all remaining
multiple-tail terms.

If $r_n(1+r_n)\ee^{2r_n}=an$ and
$S_n=\exp(nr_n(1+2r_n)/(1+r_n))/\sqrt{2r_n^2+4r_n+1}$, the result is
$v_n\sim-S_n\exp(-(1+1/a)r_n)$. More generally, arbitrary real changes to
finitely many weights and slopes preserve the equivalent, with
$1+1/a$ replaced by $(\lambda_1+w_2)/a$. Thus eventual negativity survives
finite signed perturbations. Although a large auxiliary core is used in the
proof, the leading inverse equivalent needs only the first two actions;
the one-action core is insufficient in the pure model.

A second saddle calculation determines the entire ordinary Borel transform:
if $s(s+1)=aR$ and $b=(\lambda_1+w_2)/a$, then
$-\B\V(R)\sim\exp(s\ee^{2s}/a-bs)/\sqrt{2s+1}$.
This is also its maximum-modulus equivalent. It locates the positive-direction
Borel obstruction at infinity, not at a finite Borel singularity.
The article includes complete estimates, exact finite identities, reproducible
checks, scope limitations, and twelve further research questions.
\end{abstract}

\noindent\textbf{Proof and novelty status.}
This is a conventional mathematical research draft, not a Lean-verified or
independently refereed paper. The principal result resolves a precisely
identified \emph{repository-local} conjecture. No claim of worldwide
publication priority or of solving unrestricted transseries summability is
made. The central estimates are proved here rather than imported from the
repository's positive-coefficient arguments.

\clearpage
\setcounter{tocdepth}{1}\tableofcontents
\clearpage

\section{The question and the change of viewpoint}
\label{sec:context}

\subsection{The exact repository target}
The ProveIt transseries collection develops formal inversion, support control,
large-order estimates, analytic realization, and quantitative asymptotics.
Its manuscript \emph{Finite-Core Universality and Sharp Large Order for
Countable Exponential-Feedback Transseries} explicitly labels
\begin{equation}
 v_n\sim-S_n\exp\bigl(-(1+1/a)r_n\bigr)
 \label{eq:target}
\end{equation}
as the ``Quadratic inverse equivalent'' conjecture
(TeX label \texttt{conj:quadratic-inverse}) \cite{repo-core}.
The manuscript proves the corresponding sharp inverse statement for
superquadratic feedback, but explains why that proof stops at the quadratic
boundary: its expansion parameter grows rather than tends to zero.
Its exact one-tail inverse-response identity suggests \eqref{eq:target},
while leaving dominance of that response unproved.

We address that missing dominance step. The point is not to compute more
terms in the inverse, nor to extrapolate a positive forward-coefficient
argument across a sign change. Instead, put $t=\U(q)$ and invert the
\emph{defining equation}:
\begin{equation}
 t=\sum_{j\ge1}w_jq^j\ee^{\lambda_jt},\qquad q=\V(t).
 \label{eq:direct}
\end{equation}
At fixed $t$, the feedback exponentials are now coefficients of a power
series in $q$. A finite part of that series can be inverted analytically.
Only the remaining large actions are expanded. This separates the
cancellation-producing background from the combinatorial tail estimate.

% ed. (2026-09-29): "tree" -> "commit"; the identifier is a commit (the batch-47 catalogue commit), not a Git tree object
The inspected repository commit was
\path{0c973d8f5e7ef4550f4df1bf486d890037fd9f14}.
The relevant \texttt{article.tex} blob was
\path{c2b6faf623003111e9d9e6bbb162a67429f3ce42}.
The bibliography and package provenance file give the exact path. The
inspection was targeted, not an exhaustive audit of all branches or incoming
research archives.

% ed. (2026-09-29): editorial note on status and on uncited or later repository material (negative-ray and regularity articles; weighted-type package)
\begin{ednote}
The claimed proof below has not been independently reviewed; a second,
independent claimed proof of the same conjecture is \cite{ed:qef} (see the
note after Theorem~\ref{thm:main}). The inspected snapshot also contained
two articles that this one does not cite. The negative-ray article
\cite{ed:nrs} asks, in its research question ``Signed coefficient
asymptotics and actual least error''
(\path{Negative_Ray_Summation_Exponential_Feedback/article.tex},
lines 1372--1381), for a signed asymptotic equivalent of the inverse
coefficients in the quadratic case; Theorem~\ref{thm:main} answers that
part for the inverse (the case $a=1$ with no exceptions), while the series
$\widehat P$ and the actual least truncation error there are not treated.
The regularity article \cite{ed:reg}, Theorem~\texttt{thm:borelintro}
(\path{Exponential_Feedback_Regularity_Classification/article.tex},
lines 288--303), proves the forward analogue of Theorem~\ref{thm:borel},
$\log\log B_{\mathcal U}(t)\sim2\sqrt t$ for $\lambda_j=j^2$, the same
rate as \eqref{eq:Bgrowth} at $a=1$. The later weighted-type package
\cite{ed:swt} (batch~48), which could not be seen from this snapshot,
proves for $\lambda_j\sim aj^2$ that the normalized roots
$(|v_n|\log(n+\ee)^{2n}/n!)^{1/n}$ have limsup $4a$ and that
$\limsup\log\log\max_{|z|=R}|B(z)|/\sqrt R=2\sqrt a$ ($4$ and $2$ for
$\lambda_j=j^2$); in the pure model Corollary~\ref{cor:gevrey} and
\eqref{eq:Bgrowth} upgrade both limsups to limits with the same values.
\end{ednote}

\subsection{Established tools and the proposed contribution}
Lagrange inversion is classical; we use its formal and parameter-dependent
versions in the form surveyed by Gessel \cite{gessel}. Transseries supply the
formal asymptotic setting, not an automatic assertion of analytic convergence
\cite{edgar}. The distinction between formal Borel transforms, analytic
continuation, and Laplace growth is standard in summability theory
\cite{sauzin}.

Borinsky's algebra of factorially divergent series provides an important
comparison \cite{borinsky}. Our coefficients are smaller than every fixed
positive exponential multiple of $n!$, even after multiplying by any power
of $n$. Consequently they have the \emph{zero} asymptotic expansion in every
fixed positive-$\alpha$ gamma scale used there. They are not excluded from
those classes; rather, those particular asymptotic invariants do not detect
the moving logarithmic scale resolved here.

The proposed contributions are the signed quadratic equivalent, its
finite-perturbation universality, the sharp two-action leading core, and the
explicit maximum-modulus equivalent for the inverse Borel transform. The
underlying inversion formula, Gaussian integral, and Cauchy estimates are
not claimed as new.

\section{Model, notation, and main theorems}
\label{sec:main}

\begin{definition}[Eventually quadratic feedback]
Fix $a>0$. Let $w_j,\lambda_j\in\R$, $j\ge1$, satisfy $w_1=1$ and
\begin{equation}
 w_j=1,\qquad \lambda_j=aj^2\qquad(j>M_0)
 \label{eq:model}
\end{equation}
for some fixed integer $M_0$. No positivity assumption is imposed on the
finitely many exceptional weights or slopes. Define $\U\in q\R[[q]]$ by
\begin{equation}
 \U(q)=\sum_{j\ge1}w_jq^j\ee^{\lambda_j\U(q)},
 \qquad \V=\U^{\langle-1\rangle}=\sum_{n\ge1}v_nt^n.
 \label{eq:modelU}
\end{equation}
Set
\begin{equation}
 c=\lambda_1+w_2,\qquad b=c/a,\qquad
 C_0=1+|\lambda_1|+|w_2|.
 \label{eq:cb}
\end{equation}
All parameters and all core sizes in this article are fixed while $n$ tends
to infinity.
\end{definition}

The right side of \eqref{eq:modelU} is locally finite in the $q$-adic
filtration. Its coefficient of $q^n$ depends only on coefficients of $\U$
of degrees less than $n$, because every summand has a factor $q^j$.
This proves existence and uniqueness recursively. Also
$\U(q)=q+cq^2+O(q^3)$, so its formal compositional inverse exists. Substitution
of $q=\V(t)$ gives \eqref{eq:direct}; conversely, the unique tangent-to-identity
formal solution of \eqref{eq:direct} is that inverse.

For $n>0$ let $r_n$ be the unique positive solution of
\begin{equation}
 r_n(1+r_n)\ee^{2r_n}=an.
 \label{eq:r}
\end{equation}
Define
\begin{equation}
 j_n=\frac{n}{1+r_n},\quad k_n=\frac{nr_n}{1+r_n},\quad
 D_n=2r_n^2+4r_n+1,
 \qquad
 S_n=\frac{\exp\!\left(\dfrac{nr_n(1+2r_n)}{1+r_n}\right)}{\sqrt{D_n}}.
 \label{eq:S}
\end{equation}
These are continuous saddle quantities; $j_n$ need not be an integer.
In particular,
$r_n\sim\tfrac12\log n$, $j_n\sim2n/\log n$, and
$k_n\sim n$.

\begin{theorem}[Signed quadratic inverse universality]
\label{thm:main}
Under \eqref{eq:model},
\begin{equation}
 \boxed{\displaystyle v_n\sim-S_n\ee^{-br_n}.}
 \label{eq:main}
\end{equation}
In particular, $v_n<0$ for all sufficiently large $n$, including when some
of the exceptional primitive weights are negative.
For the pure model $w_j=1$, $\lambda_j=aj^2$ for every $j$,
this is precisely \eqref{eq:target}.
\end{theorem}

The theorem concerns formal coefficients. For $t>0$ and $q\ne0$, the
quadratic tail in \eqref{eq:direct} does not even tend to zero in modulus.
Thus no claim about an ordinary positive-real inverse obtained by summing
that kernel is implicit in the theorem.

% ed. (2026-09-29): editorial note naming the other, independent claimed proof (batch 49) and the checked agreement
\begin{ednote}
A second, independent claimed proof of \eqref{eq:target} is
\cite{ed:qef}, Theorem~\texttt{thm:main}, written from the same snapshot
without knowledge of this article (neither cites the other; neither has
been independently reviewed). It inverts the kernel first, groups the
exact tree expansion by large actions and merges multiple large actions in
a positive majorant. Its hypotheses differ and neither set contains the
other: it requires nonnegative changes of finitely many weights and
slopes, but allows eventual slopes $aj^2+dj+e$, with factor
$\exp((d-\lambda_1-w_2)r_n/a)$, which is \eqref{eq:main} at $d=0$. Its
exact inverse coefficients equal the ones recorded here at every common
degree (through $220$ for $a=1$ and $180$ for $a=2$), and its positive-ray
Borel equivalent is the one in Theorem~\ref{thm:borel}. It also proves the
sharp relative error of every fixed core with at least two actions and
the constant $\ee^{\kappa/(2a)}$ for a slope perturbation
$\kappa j/\log j$; the signed exceptions, the one-action criterion of
Theorem~\ref{thm:two} and the maximum-modulus equivalent are proved only
here.
\end{ednote}

For a fixed $M\ge1$, define the analytic core and its inverse by
\begin{equation}
 H_M(q,t)=\sum_{j=1}^M w_jq^j\ee^{\lambda_jt},\qquad
 H_M(Q_M(t),t)=t,
 \qquad
 A_M(t)=\frac{1}{\partial_qH_M(Q_M(t),t)}.
 \label{eq:core}
\end{equation}
The local implicit function theorem applies because
$\partial_qH_M(0,0)=1$.
The formal one-tail response is
\begin{equation}
 L_M(t)=A_M(t)\sum_{j>M}w_jQ_M(t)^j\ee^{\lambda_jt},
 \qquad \ell_{n,M}=[t^n]L_M(t).
 \label{eq:L}
\end{equation}
This is a locally finite formal series, not a convergent sum at positive $t$.

\begin{theorem}[The leading core is exactly two actions]
\label{thm:two}
Under \eqref{eq:model}, every fixed $M\ge2$ satisfies
\begin{equation}
 \ell_{n,M}\sim S_n\ee^{-br_n},\qquad
 v_n\sim-\ell_{n,M}.
 \label{eq:two}
\end{equation}
For the one-action core,
\begin{equation}
 \ell_{n,1}\sim S_n\ee^{-(\lambda_1/a)r_n},\qquad
 \frac{v_n}{-\ell_{n,1}}\sim\ee^{-(w_2/a)r_n}.
 \label{eq:one}
\end{equation}
Hence $M=2$ is the smallest sufficient core in the pure model.
For a general finite perturbation, $M=1$ is sufficient exactly when $w_2=0$.
\end{theorem}

The proof uses a larger \emph{auxiliary} core to control the remainder.
Theorem~\ref{thm:two} is not obtained by assuming that a direct absolute bound
works optimally at $M=2$.

\begin{theorem}[Sharp entire Borel growth]
\label{thm:borel}
The ordinary Borel transform
\begin{equation}
 B(z)=\B\V(z):=\sum_{n\ge1}\frac{v_n}{n!}z^n
 \label{eq:Borel}
\end{equation}
is entire. For $R>0$ put
\begin{equation}
 s=s(R)=\frac{\sqrt{1+4aR}-1}{2}.
 \label{eq:sR}
\end{equation}
Then, as $R\to+\infty$,
\begin{equation}
 \boxed{\displaystyle
 -B(R)\sim \max_{|z|=R}|B(z)|
 \sim \frac{\exp\!\left(\dfrac{s\ee^{2s}}a-bs\right)}{\sqrt{2s+1}}.}
 \label{eq:Bmain}
\end{equation}
In particular,
\begin{equation}
 \lim_{R\to\infty}\frac{\log\log\max_{|z|=R}|B(z)|}{\sqrt R}=2\sqrt a,
 \label{eq:Bgrowth}
\end{equation}
and the ordinary positive-ray Borel--Laplace integral diverges for every
positive Laplace parameter.
\end{theorem}

\section{Exact formal inversion with a frozen parameter}
\label{sec:formal}

\subsection{A marked tail expansion}
Fix $M\ge\max(M_0,2)$ temporarily. Let $P_M(s,t)$ denote the analytic inverse
of $q\mapsto H_M(q,t)$, so $H_M(P_M(s,t),t)=s$.
The first variable $s$ and the external parameter $t$ must be kept separate.
In particular, differentiation below is with respect to $s$, with $t$
\emph{frozen}.

Introduce a marker $z$ for primitive actions larger than $M$ and let $V_z(t)$
solve
\begin{equation}
 H_M(V_z(t),t)+z\sum_{j>M}V_z(t)^j\ee^{aj^2t}=t.
 \label{eq:marked}
\end{equation}
For ordered tuples $\boldsymbol{h}=(h_1,\ldots,h_m)$ with $h_i>M$, write
\begin{equation}
 J=\sum_i h_i,\qquad
 \Lambda=a\sum_i h_i^2,\qquad
 \mathcal D_{m,J}(t)=\frac1{m!}
 \left.\partial_s^{m-1}\!\left(P_{M,s}(s,t)P_M(s,t)^J\right)\right|_{s=t}.
 \label{eq:Dmj}
\end{equation}

\begin{lemma}[Exact core-resummed expansion]
\label{lem:formal}
Coefficientwise in $t$,
\begin{equation}
 V_z(t)=Q_M(t)+
 \sum_{m\ge1}(-z)^m
 \sum_{h_1,\ldots,h_m>M}\ee^{\Lambda t}\mathcal D_{m,J}(t).
 \label{eq:exact-marked}
\end{equation}
The contribution with $m=1$ is $-zL_M(t)$.
Moreover, $\mathcal D_{m,J}(t)$ is divisible by $t^{J-m+1}$.
\end{lemma}
\begin{proof}
Put $s=H_M(V_z(t),t)$ in \eqref{eq:marked}. Then
$s=t-zT(P_M(s,t),t)$, with $T(q,t)=\sum_{j>M}q^j\ee^{aj^2t}$.
Lagrange's expansion around $s=t$, applied to $P_M(s,t)$, gives
\[
 P_M(s,t)=P_M(t,t)+
 \sum_{m\ge1}\frac{(-z)^m}{m!}
 \left.\partial_s^{m-1}\bigl(P_{M,s}(s,t)T(P_M(s,t),t)^m\bigr)\right|_{s=t}.
\]
Expanding the $m$th power gives \eqref{eq:exact-marked}.
For a completely formal justification, truncate the primitive action index
and the desired $t$-degree first, perform the finite formal calculation,
and then remove the truncation. Since $P_M(s,t)$ is divisible by $s$,
the differentiated product is divisible by $s^{J-m+1}$.
After $s=t$, a tuple can contribute to $[t^n]$ only if
$J-m+1\le n$. With $h_i\ge M+1$, this implies $mM+1\le n$ and bounds
every $h_i$. Thus all relevant sums are finite. Finally
$P_{M,s}(t,t)=A_M(t)$, proving the $m=1$ assertion.
\end{proof}

This formula is where the inverse cancellation is handled. The coefficients
of $P_M$, and therefore of $\mathcal D_{m,J}$, may have either sign. We do
not replace the core by the absolute values of its Taylor coefficients.
We estimate its analytic values on a shrinking circle only after it has
been inverted exactly.

\subsection{An independent finite identity in the pure model}
The pure model admits a useful separate check. Set $q=\ee^{-at}z$ in
\eqref{eq:direct}. Then
\begin{equation}
 t=z+\sum_{d\ge1}z^{d+1}\ee^{a d(d+1)t}.
 \label{eq:excesskernel}
\end{equation}
Inverting in $z$ by Lagrange, while keeping the $t$ in the exponentials
fixed, yields
\begin{equation}
\begin{split}
 v_n={}&\frac{(-a)^{n-1}}{(n-1)!}\\
 &+\sum_{D=1}^{n-1}\ \sum_{\substack{m\ge1,\ d_1+\cdots+d_m=D\\ d_i\ge1}}
 \frac{(-1)^m(D+m)!}{m!(D+1)!}
 \frac{\left[a\left(\sum_i d_i(d_i+1)-1\right)\right]^{n-D-1}}
 {(n-D-1)!}.
\end{split}
\label{eq:explicit}
\end{equation}
Indeed the $m$th perturbation contributes
$(-1)^m\partial_z^{m-1}(z^{D+m})/m!$ at $z=t$, followed by multiplication
by $\ee^{-at}$. Formula~\eqref{eq:explicit} is finite and exact, but it is
not a stable numerical large-order algorithm: substantial signed
cancellation remains. Its role in the package is independent low-order
verification of the direct triangular recurrence.

\section{The bare action saddle}
\label{sec:bare}

For integers $1\le j\le n$, put $k=n-j$ and
\begin{equation}
 T_{n,j}=\frac{(aj^2)^k}{k!},\qquad
 Z_n=\sum_{j=1}^nT_{n,j}.
 \label{eq:T}
\end{equation}
At $k=0$ the expression is $1$. These positive quantities are comparison
weights, not inverse coefficients.

\begin{lemma}[Bare saddle and slowly varying tilt]
\label{lem:bare}
For every fixed real $\beta$,
\begin{equation}
 Z_n\sim S_n,\qquad
 \sum_{j=1}^nT_{n,j}\exp\!\left(-\beta\frac{n-j}{j}\right)
 \sim S_n\ee^{-\beta r_n}.
 \label{eq:bare-tilt}
\end{equation}
The relevant action window is centered at $j_n$, with variance
\begin{equation}
 \sigma_n^2=\frac{k_n}{D_n}\sim\frac{2n}{(\log n)^2}.
 \label{eq:sigma}
\end{equation}
For any fixed $C>0$, comparison weights multiplied by
$\sqrt{n+1}\exp(Cj+Cn/j)$ have exponentially negligible total mass outside
\begin{equation}
 \I_n=\left[\frac{n}{2\log n},\frac{8n}{\log n}\right]
 \label{eq:coarse}
\end{equation}
when their coefficient radii $k/(aj^2)$ are bounded by a sufficiently small
fixed constant. Indices for which those radii are larger can instead be
bounded at a fixed radius by $\exp(O(n))$.
\end{lemma}
\begin{proof}
Stirling's formula, away from $k=0$, gives
\begin{equation}
 T_{n,j}=\frac{\ee^{f_n(j)}}{\sqrt{2\pi(n-j)}}
 \left(1+O\bigl((n-j)^{-1}\bigr)\right),\quad
 f_n(j)=(n-j)\left(1+\log\frac{aj^2}{n-j}\right).
 \label{eq:f}
\end{equation}
Differentiation gives
\begin{align}
 f_n'(j)&=-\log\frac{aj^2}{n-j}+\frac{2(n-j)}j,\label{eq:fd1}\\
 f_n''(j)&=-\frac4j-\frac{2(n-j)}{j^2}-\frac1{n-j}<0.\label{eq:fd2}
\end{align}
Thus there is a unique stationary point. Setting $r=(n-j)/j$ in
$f_n'(j)=0$ gives \eqref{eq:r}, hence $j=j_n$. At that point,
\begin{equation}
 f_n(j_n)=\frac{nr_n(1+2r_n)}{1+r_n},\qquad
 -f_n''(j_n)=\frac{D_n}{k_n}.
 \label{eq:fstationary}
\end{equation}
On $j\asymp n/\log n$, the third derivative is
$O(n/j^3)$ and the fourth is $O(n/j^4)$. Therefore, for
$j=j_n+\sigma_n x$ with $|x|\le\log n$,
\[
 f_n(j)=f_n(j_n)-x^2/2+O((1+|x|)^3/\sqrt n).
\]
The Stirling prefactor is uniform there. Summation of the lattice Gaussian,
whose mesh is $1/\sigma_n\to0$, gives
\[
 \sum_jT_{n,j}\sim
 \frac{\ee^{f_n(j_n)}}{\sqrt{2\pi k_n}}\sqrt{2\pi}\sigma_n=S_n.
\]
Strict concavity and \eqref{eq:fd2} bound the portion of $\I_n$ outside that
window by a polynomial factor times $\ee^{-c(\log n)^2}S_n$.

For the tilt, put $r(j)=n/j-1$. In the same window,
\[
 |r(j)-r_n|\le C\frac{n}{j_n^2}\sigma_n\log n
 =O((\log n)^2/\sqrt n)=o(1).
\]
Thus the tilt is uniformly $\ee^{-\beta r_n}(1+o(1))$ in the central window.
On $\I_n$ it is at most a fixed power of $n$, so the Gaussian tail bound
remains sufficient.

Here are details of the coarser localization, also needed below. Write
$x=j\log n/n$. Uniformly for $x$ in a fixed compact subinterval of $(0,\infty)$,
\begin{equation}
 \frac{f_n(j)}n=
 \log n-2\log\log n+\log a+1+2\log x-x+o(1).
 \label{eq:macro}
\end{equation}
The function $\phi(x)=2\log x-x$ has its unique maximum at $x=2$ and a
strictly smaller value on $x\le1/2$ or $x\ge8$. To make the off-window
comparison uniform, put $\theta=j/n$ and $L=\log n$. There is the exact identity
\begin{align}
 \frac{f_n(j)}n&=L-2\log L+\log a+1+\phi(x)+E_a(\theta),\label{eq:global-macro}\\
 E_a(\theta)&=-2\theta\log\theta-\theta(\log a+1)
                -(1-\theta)\log(1-\theta).\nonumber
\end{align}
With $0\log0=0$, $E_a$ is bounded on $[0,1]$ and tends to zero as
$\theta\downarrow0$. In the small-radius range $k/(aj^2)\le\delta$,
either $j\ge n/2$ or $j\ge\sqrt{n/(2a\delta)}$. Thus the decoration,
divided by $n$, is $C\theta+C/j+o(1)$, with $C/j=o(1)$ uniformly.
For $x\le1/2$ and for $8\le x\le K$, where $K$ is fixed, both
$E_a(\theta)$ and $C\theta$ tend uniformly to zero. The strict loss of
$\phi$ proves the desired bound there. Choose $K>8$ so large that
$\phi(K)+\sup E_a+C<\phi(2)-1$. As $\phi$ is decreasing after $2$,
\eqref{eq:global-macro} gives a loss at least $n/2$ for $x\ge K$,
for large $n$. Together these estimates give a loss at least $\eta n$
outside \eqref{eq:coarse}. The Stirling prefactor costs only a polynomial;
the endpoint $k=0$ is handled directly and contributes an exponential
decoration, still negligible compared with $S_n$.

Finally, if $k/(aj^2)>\delta$ for fixed $\delta>0$, then $j=O(\sqrt n)$.
A Cauchy bound at radius $\delta$ has logarithm at most
$n\log(1/\delta)+aj^2\delta+O(j)=O(n)$.
This is negligible compared with
$\log S_n=n\log n-2n\log\log n+O(n)$.
This proves both the localization statements and \eqref{eq:bare-tilt}.
\end{proof}

The only use of the last, decorated assertion will have a fixed $C$.
No uniformity as the core or the model parameters vary with $n$ is being
asserted.

\section{A signed finite-core reduction with arbitrary algebraic accuracy}
\label{sec:reduction}

\subsection{Core bounds independent of the number of retained actions}
The first jet of $P_M(s,t)/s$ is independent of $M$ once $M\ge2$:
\begin{equation}
 \log\frac{P_M(s,t)}s=-\lambda_1t-w_2s+O((|s|+|t|)^2).
 \label{eq:Pjet}
\end{equation}
For each fixed $M$ there is consequently a radius $\delta_M>0$ on which
\begin{equation}
 |P_M(s,t)|\le |s|\exp(C_0(|s|+|t|)),\qquad |P_{M,s}(s,t)|\le2.
 \label{eq:Pbounds}
\end{equation}
The radius may depend on all the retained data, but the displayed constant
$C_0=1+|\lambda_1|+|w_2|$ does not grow with $M$. This distinction permits
us to choose a larger auxiliary core without a circular dependence of the
power-saving exponent on that choice.

\begin{lemma}[Derivative estimate]
\label{lem:D-bound}
For every fixed sufficiently large $M$, every tuple in
\eqref{eq:Dmj}, and every sufficiently small $\tau>0$,
\begin{equation}
 \sup_{|t|=\tau}|\mathcal D_{m,J}(t)|
 \le \frac{4^mJ^{m-1}}{m!}\,
 \tau^{J-m+1}\exp(3C_0J\tau).
 \label{eq:D-bound}
\end{equation}
The smallness threshold for $\tau$ depends on $M$ but not on the tuple.
\end{lemma}
\begin{proof}
For $m=1$, this follows directly from \eqref{eq:Pbounds}.
For $m\ge2$, take a Cauchy circle in the $s$ variable centered at $s=t$
with radius $\rho=(m-1)\tau/J$. Since $J\ge m(M+1)$, this circle lies in
$|s|\le2\tau$. On it,
\[
 |P_M(s,t)|^J\le
 \tau^J\left(1+\frac{m-1}J\right)^J\ee^{3C_0J\tau}
 \le \tau^J\ee^{m-1+3C_0J\tau}.
\]
Cauchy's derivative estimate and the factor $1/m!$ in
\eqref{eq:Dmj} give
\[
 |\mathcal D_{m,J}|
 \le\frac2m\left(\frac{\ee J}{m-1}\right)^{m-1}
 \tau^{J-m+1}\ee^{3C_0J\tau}.
\]
Since $(m-1)!\le(m-1)^{m-1}$ and $2\ee^{m-1}\le4^m$, this is
\eqref{eq:D-bound}.
\end{proof}

\subsection{Merging the tail actions}
For a tuple let
\begin{equation}
 d_i=h_i-1\ge M,\quad D=\sum_i d_i,\quad j=D+1=J-m+1,\quad k=n-j.
 \label{eq:merge-notation}
\end{equation}
The effective single action $j$ has the same minimal $t$-degree as the tuple.
Its quadratic excess over the tuple is
\begin{equation}
 \Delta=j^2-\sum_i h_i^2
 =D^2-\sum_i d_i^2-(m-1)
 =2\sum_{i<l}d_id_l-(m-1)\ge0.
 \label{eq:defect}
\end{equation}
This exact identity is the tail-suppression mechanism.

\begin{lemma}[Coefficient bound by a merged action]
\label{lem:merged}
For $j\in\I_n$, $m\ge1$, and all sufficiently large $n$,
\begin{equation}
 \left|[t^n]\ee^{\Lambda t}\mathcal D_{m,J}(t)\right|
 \le C\sqrt{n+1}\,T_{n,j}\,
 \frac{4(8j)^{m-1}}{m!}\,
 \exp(6C_0j\tau)\exp(-a\tau\Delta),
 \qquad \tau=\frac{k}{aj^2}.
 \label{eq:merged-bound}
\end{equation}
Here and below constants can depend on the fixed model and core.
\end{lemma}
\begin{proof}
In \eqref{eq:D-bound}, $J=j+m-1\le2j$. Cauchy's coefficient bound at
$|t|=\tau$ is therefore at most
\[
 \frac{4(8j)^{m-1}}{m!}\tau^{-k}
 \ee^{a(j^2-\Delta)\tau+6C_0j\tau}.
\]
At $\tau=k/(aj^2)$,
\[
 \frac{\tau^{-k}\ee^k}{T_{n,j}}=
 \frac{k!\ee^k}{k^k}\le C\sqrt{k+1}
\]
by Stirling's bound. On $\I_n$, $\tau\to0$, so the analytic estimates
apply uniformly for large $n$.
\end{proof}

\subsection{One giant tail action versus several}
We now sum the bounds without assuming signs of any core coefficients.
Let $L=\max_i d_i$ and $E=D-L$.
If $L\ge D/2$ and $m\ge2$, then $E\ge M$ and
\begin{equation}
 \Delta\ge2LE-(m-1)\ge(D-1/M)E\ge jE/2
 \label{eq:giantdefect}
\end{equation}
for large $j$. On $\I_n$, $k/j\ge(\log n)/9$ eventually. Hence
\begin{equation}
 \ee^{-a\tau\Delta}\le n^{-E/20}.
 \label{eq:powertail}
\end{equation}
The constants $9$ and $20$ are deliberately loose.

For fixed $m$, choose a position for the largest component (at most $m$
choices) and sum the other $m-1$ components. Dropping the upper bound on
their total only enlarges the bound. Put
\[
 z_n=\sum_{d\ge M}n^{-d/20}
 =\frac{n^{-M/20}}{1-n^{-1/20}}\le2n^{-M/20}
\]
for sufficiently large $n$. The sum of the tuple factors in
\eqref{eq:merged-bound} with a giant and $m\ge2$ is at most
\begin{equation}
 4\sum_{m\ge2}\frac{(8jz_n)^{m-1}}{(m-1)!}
 =4(\ee^{8jz_n}-1)=O(jn^{-M/20})
 \label{eq:giantsum}
\end{equation}
when $M>20$.

If no component is at least $D/2$, then
$\sum_i d_i^2\le LD<D^2/2$ and $m-1\le D/M$.
Thus $\Delta\ge j^2/3$ for large $j$. The suppression in
\eqref{eq:merged-bound} is at most $\ee^{-k/3}$.
There are at most $2^D$ ordered compositions of $D$, and
\[
 \sum_{m\ge1}\frac{4(8j)^{m-1}}{m!}\le4\ee^{8j}.
\]
The total tuple factor in the no-giant case is consequently at most
$4\ee^{9j-k/3}=\ee^{-\eta n}$ on $\I_n$ for large $n$.
All these estimates allow overcounting ties for the largest component.

Finally, on $\I_n$,
\begin{equation}
 \ee^{6C_0j\tau}
 =\exp\!\left(\frac{6C_0}a\frac{k}j\right)
 \le n^{12C_0/a}.
 \label{eq:corepoly}
\end{equation}
Combining \eqref{eq:merged-bound}--\eqref{eq:corepoly}, summing over $j$,
and using Lemma~\ref{lem:bare}, bounds the multi-tail contribution in the
coarse window by
\begin{equation}
 C S_n\,n^{3/2+12C_0/a-M/20}+S_n\ee^{-\eta n}.
 \label{eq:remainder-coarse}
\end{equation}

\subsection{Off-window contributions and the finite-core theorem}
Outside $\I_n$, first ignore the nonnegative merger defect. The sum of all
tuple combinatorial factors with fixed $j$ is bounded by $\ee^{Cj}$, by the
same composition estimate just used. In the small-radius range the total
coefficient majorant is therefore
\[
 C\sqrt{n+1}\,T_{n,j}\exp(Cj+Cn/j).
\]
Lemma~\ref{lem:bare} makes its off-window sum exponentially negligible
relative to $S_n$. In the complementary range $j=O(\sqrt n)$, use a fixed
Cauchy radius small enough for the core instead. The logarithm of the total
bound is $O(n)$, again negligible. The coefficient of the analytic series
$Q_M$ is also at most exponential in $n$.

\begin{theorem}[Arbitrary algebraic accuracy after a finite core]
\label{thm:reduction}
For every fixed $A>0$ there is a fixed $M$ such that
\begin{equation}
 \boxed{\displaystyle
 v_n+\ell_{n,M}=O\!\left(S_n n^{-A-|b|}\right).}
 \label{eq:finite-accuracy}
\end{equation}
One sufficient, non-optimal choice is any integer $M\ge\max(M_0,2)$ satisfying
\begin{equation}
 \frac M{20}>A+3+\frac{12C_0}{a}+|b|.
 \label{eq:corechoice}
\end{equation}
\end{theorem}
\begin{proof}
Use the exact expansion \eqref{eq:exact-marked} at $z=1$.
The $m=1$ term is $-L_M$. The remaining tuple terms satisfy
\eqref{eq:remainder-coarse} in $\I_n$ and the exponentially smaller
estimates outside it. The strict inequality in \eqref{eq:corechoice}
absorbs the factor $n^{3/2}$ and leaves more than the required saving.
The analytic core coefficient is negligible on the same scale.
\end{proof}

\begin{remark}[Why this proves a signed statement]
The forward coefficients never enter the proof. The inverse response itself
has size $S_n$ times a possibly small power of $n$; a merely $o(S_n)$ error
would not suffice. The explicit saving in \eqref{eq:finite-accuracy} beats
that polynomial suppression. This is exactly the obstacle avoided by using
a sufficiently large auxiliary core before reducing to the minimal core.
\end{remark}

\section{Evaluation of the one-tail inverse response}
\label{sec:one-tail}

For every fixed $M\ge2$, the diagonal core has expansion
\begin{equation}
 Q_M(t)=t-ct^2+O(t^3),\qquad A_M(t)=1+O(t).
 \label{eq:Qjet}
\end{equation}
Consequently, in a fixed disk about zero,
\begin{equation}
 \log\frac{Q_M(t)}t=-ct+R_M(t),\qquad R_M(t)=O(t^2).
 \label{eq:logQ}
\end{equation}
For large tail actions the one-tail summand is
\begin{equation}
 A_M(t)Q_M(t)^j\ee^{aj^2t}
 =t^j\ee^{\alpha_jt}E_j(t),\qquad
 \alpha_j=aj^2-cj,
 \quad E_j(t)=A_M(t)\ee^{jR_M(t)}.
 \label{eq:factor-one}
\end{equation}
For all sufficiently large $j$, $\alpha_j>0$, even when $c$ is negative.

\begin{lemma}[Uniform inner coefficient estimate]
\label{lem:inner}
Uniformly for integers $j\in\I_n$, with $k=n-j$,
\begin{equation}
 [t^k]\ee^{\alpha_jt}E_j(t)
 =\frac{\alpha_j^k}{k!}
 \left(1+O\left(\frac{(\log n)^3}{n}\right)\right).
 \label{eq:inner}
\end{equation}
\end{lemma}
\begin{proof}
Use the radius $\rho=k/\alpha_j$. On $\I_n$,
$\rho=O((\log n)^2/n)$ and $j\rho^2=O((\log n)^3/n)$.
By analyticity and \eqref{eq:logQ},
\[
 \sup_{|t|=\rho}|E_j(t)-1|\le C(\rho+j\rho^2)=:\varepsilon_n.
\]
The Cauchy integral for the error is at most
\[
 \varepsilon_n\rho^{-k}\frac1{2\pi}
 \int_{-\pi}^{\pi}\ee^{k\cos\theta}\,\dd\theta
 \le C\varepsilon_n\rho^{-k}\frac{\ee^k}{\sqrt{k+1}}.
\]
The last bound follows from $1-\cos\theta\ge2\theta^2/\pi^2$ on
$[-\pi,\pi]$. Stirling's lower bound shows that this is at most
$C\varepsilon_n\alpha_j^k/k!$. The coefficient with $E_j=1$ is exactly
$\alpha_j^k/k!$. Notice that the entire factor $\ee^{-cjt}$ was absorbed
into $\alpha_j$ before taking absolute values; failing to do this would
lose the polynomial cancellation we need.
\end{proof}

Since $\alpha_j/(aj^2)=1-b/j$, on $\I_n$ we have
\begin{equation}
 \frac{\alpha_j^k}{(aj^2)^k}
 =\exp\left(-b\frac{k}j+O(k/j^2)\right)
 =\ee^{-bk/j}\left(1+O((\log n)^2/n)\right).
 \label{eq:tilt}
\end{equation}
Summing \eqref{eq:inner} and using Lemma~\ref{lem:bare} gives
\begin{equation}
 \ell_{n,M}\sim S_n\ee^{-br_n}.
 \label{eq:L-asymptotic}
\end{equation}
The off-window terms are controlled by the $m=1$ version of the analytic
majorant used in Section~\ref{sec:reduction}. Finitely many nonquadratic
or exceptional summands are analytic at zero and contribute only
exponentially growing coefficients. Thus \eqref{eq:L-asymptotic} holds for
\emph{every} fixed $M\ge2$, even if that core does not contain all exceptional
actions.

\begin{proof}[Proof of Theorems~\ref{thm:main} and \ref{thm:two}]
Choose the auxiliary $M$ from Theorem~\ref{thm:reduction}, for instance with
$A=1$. Since $r_n<\log n$ eventually,
$\ee^{-br_n}\ge n^{-|b|}$. Hence the error in
\eqref{eq:finite-accuracy} is $o(S_n\ee^{-br_n})$.
Combine it with \eqref{eq:L-asymptotic} to obtain \eqref{eq:main}.
The same one-tail equivalent for any other fixed $M\ge2$ then proves
\eqref{eq:two} without a second remainder estimate at that smaller core.

For $M=1$, $Q_1(t)=t\ee^{-\lambda_1t}$ and
$A_1(t)=\ee^{-\lambda_1t}$. The preceding one-tail calculation applies with
$c=\lambda_1$; the additional factor $A_1(t)=1+O(t)$ does not alter the
equivalent. Division by \eqref{eq:main} gives \eqref{eq:one}.
As $r_n\to\infty$, its right side tends to $1$ exactly for $w_2=0$,
to $0$ for $w_2>0$, and to $+\infty$ for $w_2<0$.
\end{proof}

\subsection{What the suppression factor measures}
At the bare inner saddle $t\approx k/(aj^2)$,
\[
 j\log(Q_M(t)/t)=-cjt+O(jt^2)
 =-\frac ca\frac kj+o(1).
\]
The outer saddle has $k/j=r_n$. This gives the entire leading penalty
$\ee^{-br_n}$. Only the quadratic coefficient of the inverse core survives.
All higher diagonal core coefficients are multiplied by quantities beginning
with $jt^2=O((\log n)^3/n)$.

This mechanism differs from simply substituting the forward coefficient
asymptotic into a few inverse-composition terms. In the pure model, the
repository's forward result is
$u_n\sim S_n\exp((1+1/a)r_n^2)$ \cite{repo-core}.
\emph{Combining that external forward theorem with our independent inverse
theorem} gives
\begin{equation}
 \frac{-v_n}{u_n}\sim
 \exp\bigl(-(1+1/a)(r_n^2+r_n)\bigr).
 \label{eq:forward-comparison}
\end{equation}
The forward theorem is needed only for this comparison, not for any of our
inverse or Borel proofs.

\subsection{An explicit two-action formula}
For the pure model,
\begin{equation}
 Q_2(t)=\frac{\ee^{-3at}}2\left(\sqrt{1+4t\ee^{2at}}-1\right),\qquad
 A_2(t)=\frac{\ee^{-at}}{\sqrt{1+4t\ee^{2at}}}.
 \label{eq:Q2}
\end{equation}
These are the branches analytic at zero. Let
$C(z)=(1-\sqrt{1-4z})/(2z)$ be the Catalan series.
Then $Q_2=t\ee^{-at}C(-t\ee^{2at})$, and the elementary identity
\[
 \frac{C(z)^j}{\sqrt{1-4z}}=
 \sum_{h\ge0}\binom{2h+j}{h}z^h
\]
(obtainable by Lagrange inversion) gives the finite coefficient formula
\begin{equation}
 \ell_{n,2}=\sum_{j=3}^n\ \sum_{h=0}^{n-j}
 (-1)^h\binom{2h+j}{h}
 \frac{\left[a(j^2-j-1+2h)\right]^{n-j-h}}{(n-j-h)!}.
 \label{eq:L2exact}
\end{equation}
Thus a completely explicit finite sum already has the correct full inverse
equivalent. It is still a signed sum; exact arithmetic or adequate working
precision is necessary.

\section{Consequences for coefficient growth and formal transseries}
\label{sec:consequences}

\begin{corollary}[Ratio law and exact Gevrey threshold]
\label{cor:gevrey}
Under \eqref{eq:model},
\begin{align}
 \frac{v_n}{v_{n-1}}&\sim\ee^{2r_n}
 \sim\frac{4an}{(\log n)^2},\label{eq:ratio}\\
 |v_n|^{1/n}&\sim\frac{4an}{\ee(\log n)^2},\label{eq:root}\\
 \log\frac{|v_n|}{n!}
 &=-2n\log\log n+n\log(4a)+o(n).\label{eq:fact-growth}
\end{align}
The inverse has radius of convergence zero. It is Gevrey-$1$ of zero
exponential type: for every $A>0$ there is $C_A<\infty$ with
$|v_n|\le C_AA^nn!$. It is not Gevrey-$s$ for any $s<1$.
\end{corollary}
\begin{proof}
Extend $r_n$ smoothly to positive real $n$. Implicit differentiation gives
\begin{equation}
 r'(n)=\frac{r(1+r)}{n(2r^2+4r+1)}.
 \label{eq:rprime}
\end{equation}
For $F(n)=nr(1+2r)/(1+r)$, the envelope identity is $F'(n)=2r$.
The logarithmic derivative of $\sqrt{2r^2+4r+1}$ is $O(1/n)$ and
$br'(n)=O(1/n)$. Thus the ratio of successive smooth equivalents is
$\ee^{2r_n}(1+o(1))$. The two relative errors in Theorem~\ref{thm:main}
also have ratio tending to one. This proves \eqref{eq:ratio}.

Using \eqref{eq:r},
\[
 \frac1n\log|v_n|=2r_n-1+\frac1{1+r_n}+o(1),\qquad
 \ee^{2r_n}=\frac{an}{r_n(1+r_n)}.
\]
This proves \eqref{eq:root}; Stirling then gives \eqref{eq:fact-growth}.
The conclusions follow directly. For $s<1$,
$\log(|v_n|/(n!)^s)=(1-s)n\log n-2n\log\log n+O(n)$,
which exceeds $n\log A$ eventually for every fixed $A$.
\end{proof}

\begin{corollary}[Flatness in every fixed positive factorial scale]
For all $\alpha>0$, all real $\beta$, and every integer $K\ge0$,
\begin{equation}
 v_n=o\bigl(\alpha^{n+\beta}\Gamma(n+\beta-K)\bigr).
 \label{eq:gammaflat}
\end{equation}
Thus the gamma-scale asymptotic derivation of \cite{borinsky}, with any
fixed positive $\alpha$, assigns the zero expansion to this divergent inverse.
\end{corollary}
\begin{proof}
The logarithm of the ratio in \eqref{eq:gammaflat} equals
$-2n\log\log n+O(n)+O(\log n)$, which tends to $-\infty$.
This is exactly the all-zero case of the fixed gamma-scale definition.
\end{proof}

\begin{corollary}[Least formal term]
For $t>0$, let $N(t)$ minimize $|v_n|t^n$ over sufficiently large indices
past the eventual-sign threshold. Then, as $t\downarrow0$,
\begin{equation}
 N(t)\sim\frac{(\log(1/t))^2}{4at}.
 \label{eq:least}
\end{equation}
\end{corollary}
\begin{proof}
The terms eventually grow without bound for fixed $t$, by
\eqref{eq:root}. Their successive ratios are asymptotic to
$t\ee^{2r_n}$, an increasing smooth comparison.
At its crossing of $1$, $r=\frac12\log(1/t)$ and the corresponding real
index is $r(1+r)/(at)$. The uniform eventual relative bounds furnished by
\eqref{eq:ratio} confine any minimizing index to a relative $o(1)$
window about this value.
\end{proof}

This is a statement about formal terms, not a remainder estimate for an
actual analytic realization. Substituting $t=\ee^{-x}$ gives an exponential
transseries with a least-term index asymptotic to $x^2\ee^x/(4a)$.
No summability conclusion is inferred from this formal optimization.

\section{Entire Borel growth: a second, explicit saddle}
\label{sec:borel}

\subsection{Passing from coefficient equivalents to an entire function}
Equation~\eqref{eq:fact-growth} proves that $B$ in \eqref{eq:Borel} is entire.
By Theorem~\ref{thm:main}, its coefficients are eventually negative.
Let
\[
 \widetilde B(R)=\sum_{n\ge1}\frac{S_n\ee^{-br_n}}{n!}R^n.
\]
All terms are positive, and their sum tends to infinity faster than every
polynomial. The coefficient equivalent implies
\begin{equation}
 -B(R)\sim\widetilde B(R).
 \label{eq:Btransfer}
\end{equation}
To see this without an unjustified interchange of limits, fix $\varepsilon>0$.
For every $n$ beyond a fixed threshold, the ratio of the actual negative
coefficient to the positive comparison coefficient is between
$1-\varepsilon$ and $1+\varepsilon$. Sum those inequalities. The finitely
many discarded terms form a polynomial, negligible relative to
$\widetilde B(R)$. Then let $\varepsilon\downarrow0$.

\subsection{Stationary point and determinant}
For continuous $n>0$, use $r=r(n)$ from \eqref{eq:r} and define
\begin{equation}
 h_R(n)=F(n)+n(1-\log n)+n\log R,
 \qquad F(n)=\frac{nr(1+2r)}{1+r}.
 \label{eq:hR}
\end{equation}
By Stirling,
\begin{equation}
 \frac{S_n\ee^{-br_n}R^n}{n!}
 =\frac{\ee^{h_R(n)-br_n}}{\sqrt{2\pi nD_n}}
 \left(1+O(n^{-1})\right).
 \label{eq:Bterm}
\end{equation}
The exact derivative identities are
\begin{align}
 h_R'(n)&=2r-\log n+\log R
 =\log\frac{aR}{r(1+r)},\label{eq:hprime}\\
 h_R''(n)&=-\frac{2r+1}{n(2r^2+4r+1)}<0.\label{eq:hsecond}
\end{align}
Thus the saddle is characterized simply by $r=s$, where $s(1+s)=aR$.
Its real index and exponent are
\begin{equation}
 n_R=R\ee^{2s}=\frac{s(1+s)}a\ee^{2s},\qquad
 h_R(n_R)=\frac{s\ee^{2s}}a.
 \label{eq:Bstationary}
\end{equation}
The variance in the coefficient index is
\begin{equation}
 \Sigma_R^2=\frac{n_R(2s^2+4s+1)}{2s+1}.
 \label{eq:Bvariance}
\end{equation}

\begin{lemma}[The outer Laplace sum]
\label{lem:Bsum}
As $R\to\infty$,
\begin{equation}
 \widetilde B(R)\sim
 \frac{\ee^{h_R(n_R)-bs}}{\sqrt{2s+1}}.
 \label{eq:Bsum}
\end{equation}
\end{lemma}
\begin{proof}
The width satisfies $\Sigma_R/n_R=O(\sqrt{s/n_R})\to0$.
In a neighborhood of $n_R$ with relative width tending to zero,
\eqref{eq:rprime} gives $r(n)-s=O(|n-n_R|/n_R)$.
The amplitude $\ee^{-br}/\sqrt{2\pi nD}$ therefore has relative change
$o(1)$ on a Gaussian window.
Differentiating \eqref{eq:hsecond} gives
$h_R'''(n)=O(1/(n_R^2s))$ there, so
$h_R'''\Sigma_R^3=O(\sqrt{s/n_R})\to0$.
Taylor expansion and lattice Gaussian summation yield
\[
 \widetilde B(R)\sim
 \frac{\ee^{h_R(n_R)-bs}}{\sqrt{2\pi n_R(2s^2+4s+1)}}
 \sqrt{2\pi}\Sigma_R
 =\frac{\ee^{h_R(n_R)-bs}}{\sqrt{2s+1}}.
\]

For completeness, the tail localization follows from the global strict
concavity \eqref{eq:hsecond}. On $[n_R/2,2n_R]$,
$-h_R''$ is comparable to $1/(n_Rs)$, giving Gaussian decay outside an
expanding central window. At the endpoints the exponent loses at least
a fixed multiple of $n_R/s$. For $n\ge2n_R$, the decreasing derivative
satisfies $h_R'(n)\le h_R'(2n_R)\le-c/s$ for large $R$; thus the remaining
sum is bounded by a geometric tail after its already negligible endpoint.
The lower tail is handled by the analogous positive derivative to the
left of $n_R/2$. The amplitudes grow or decay only by fixed powers of $n$
and logarithms there, so they do not offset these exponential losses.
This also justifies discarding finitely many small indices and using
Stirling uniformly in the contributing range.
\end{proof}

\begin{proof}[Proof of Theorem~\ref{thm:borel}]
Combine \eqref{eq:Btransfer}, Lemma~\ref{lem:Bsum}, and
\eqref{eq:Bstationary}. Because $v_n<0$ eventually,
\[
 |B(R)|\le\max_{|z|=R}|B(z)|
 \le\sum_{n\ge1}|v_n|R^n/n!=-B(R)+P(R),
\]
where $P$ is a fixed polynomial accounting for the finitely many coefficients
with the other sign. Since $-B(R)$ has the growth already proved,
$P(R)=o(-B(R))$, which gives the maximum-modulus equivalent.

Taking logarithms twice in \eqref{eq:Bmain}, the leading expression is
$2s+\log(s/a)+o(1)$. As $s\sim\sqrt{aR}$, this proves
\eqref{eq:Bgrowth}. In our ordinary Borel convention the positive-ray
Laplace expression would be
\[
 \frac1t\int_0^\infty\ee^{-R/t}B(R)\,\dd R,\qquad t>0.
\]
The integrand is eventually negative, and
$\log(-B(R))-R/t\to+\infty$. It therefore does not tend to zero and its
integral diverges. The same obstruction rules out positive-direction
ordinary $1$-summability; analyticity alone is not enough without the
required exponential-growth bound \cite{sauzin}.
\end{proof}

\subsection{Interpretation and limits}
The Borel transform has no finite singularities at all. Nevertheless it
cannot be integrated by the ordinary positive-direction Laplace kernel.
The obstruction is its growth at infinity. This supplies a concrete example
where exact Gevrey-$1$ regularity, an entire Borel transform, and failure of
ordinary directional Borel summation coexist.

The formula does not classify other complex directions. Cancellations along
nonreal rays can substantially change growth. Nor does it identify a unique
acceleration procedure or a canonical sectorial realization. Those are
separate analytic problems, listed below.

\section{Exact computations and diagnostic evidence}
\label{sec:computation}

\subsection{Independent algebraic routes}
The package computes the inverse through degree $220$ in four fixed models,
using integer exponential normalization
$V_n=n!v_n$. For integer weights and slopes these are integers.
Let $P_{j,n}=n![t^n]\V(t)^j$. The identities
\begin{align}
 P_{j,n}&=\sum_{k=1}^{n-j+1}\binom nk V_kP_{j-1,n-k}
 \qquad(j\ge2),\label{eq:recP}\\
 V_n&=-\sum_{j=1}^n w_j
 \sum_{k=j}^{n-\mathbf1_{j=1}}
 \binom nk\lambda_j^{n-k}P_{j,k}
 \qquad(n\ge2),\label{eq:recV}
\end{align}
with $V_1=P_{1,1}=1$, form a triangular exact algorithm.
Formula~\eqref{eq:recV} is just the coefficient of \eqref{eq:direct}, with
the unique unknown $V_n$ isolated.

The first independent check uses \eqref{eq:explicit} through degree $12$
for $a=1$ and $a=2$ (24 assertions). The second constructs the forward
coefficients independently by the positive Lagrange identity
\begin{equation}
 u_n=\sum_{\substack{j_1+\cdots+j_m=n\\j_i\ge1,\ m\ge1}}
 \frac{(a\sum_i j_i^2)^{m-1}}{m!}
 \label{eq:forward-exact}
\end{equation}
and checks $\U(\V(t))=t$ through degree $11$ in both models (two block
assertions). A separate symbolic script verifies six identities: the action
curvature, the coefficient envelope derivative, the Borel curvature and
stationary value, and two core-jet formulas. All 32 recorded assertions pass.
These are finite exact checks, not formal proofs of the limit theorems.

\subsection{What the numerical ratios show}
Table~\ref{tab:pure} uses exact coefficients followed by 80-digit evaluation
of the smooth comparison. The errors shown are not interval certificates.
The two-core response is evaluated independently using \eqref{eq:L2exact}.

\begin{table}[htbp]\centering\small
\caption{Pure quadratic feedback. The predicted limit in both ratio columns is $1$.}
\label{tab:pure}
\begin{tabular}{rrrrr}\toprule
$a$&$n$&$r_n$&$-v_n/(S_n\ee^{-(1+1/a)r_n})$&$-v_n/\ell_{n,2}$\\\midrule
1&40&1.298005&0.607745&0.855019\\
1&80&1.519691&0.725965&0.908128\\
1&120&1.654076&0.780419&0.928350\\
1&160&1.751348&0.813338&0.939697\\
1&220&1.860772&0.844814&0.950076\\\addlinespace
2&40&1.519691&0.656706&0.930063\\
2&80&1.751348&0.768248&0.961094\\
2&120&1.890976&0.817596&0.971321\\
2&160&1.991732&0.846760&0.976658\\
2&220&2.104802&0.874148&0.981282\\\bottomrule
\end{tabular}
\end{table}

The main equivalent has substantial finite-$n$ corrections even at degree
$220$; the table must not be described as near-perfect asymptotic agreement.
The exact two-core expression absorbs more of those corrections and is
noticeably closer. Neither observation replaces the proof of tail dominance.

Two additional tests change only finitely many primitive data. In the first,
$a=1$, $w_2=-2$, $w_3=3$, $\lambda_1=0$, and $\lambda_2=7$, with the other
data pure. Here $c=-2$: the leading core factor is an enhancement
$\ee^{2r_n}$ rather than a suppression. The ratio to the predicted equivalent
is about $2.693$ at $n=220$, so the asymptotic regime is visibly slow.
In the second, $a=1$, $w_2=-1$ and all other data are pure, so $c=0$ and
the leading penalty cancels. The exact data for both models are included in
\texttt{data/diagnostics.csv}; they are stress tests of the signed-core
extension rather than evidence for a uniform onset bound.

\subsection{Reproduction and numerical safeguards}
Run
\begin{verbatim}
python code/verify.py --order 220
python code/symbolic_checks.py
pdflatex -interaction=nonstopmode -halt-on-error article.tex
pdflatex -interaction=nonstopmode -halt-on-error article.tex
\end{verbatim}
The exact coefficient algorithm needs only Python's standard library.
\texttt{mpmath} is used for the comparison values, and \texttt{sympy} for the
six symbolic checks. The scripts validate the model normalization and retain
all coefficient signs. The ordinary coefficients are recovered by dividing
stored integers by $n!$; confusing the two normalizations would invalidate
both the ratios and the Borel transform.

The full exact algorithm is cubic in the truncation order in arithmetic
operation count and uses a quadratic table. Its bit cost is larger because
coefficient sizes increase rapidly. No complexity claim suppresses that
integer-arithmetic cost.

\section{Further research questions and topics}
\label{sec:questions}

\begin{question}[Effective onset and small-$a$ uniformity]
Replace the asymptotic neighborhoods in the core estimates by computable
radii and constants. The target is a certified pair of bounds for
$-v_n/(S_n\ee^{-br_n})$ with an explicit onset depending on $a$ and the
exceptional data. The proof here is not uniform as $a\downarrow0$; the
factor $b=1+1/a$ makes that limit especially delicate.
\end{question}

\begin{question}[An optimal signed remainder at the two-action core]
Theorem~\ref{thm:two} proves $v_n\sim-\ell_{n,2}$ but obtains it through a
large auxiliary core. Determine the first term of
$v_n+\ell_{n,2}$ and a sharp relative error. The natural next contribution
is one additional primitive action $3$, dressed by the two-action core;
its cancellation and combinatorial factors should be kept exact.
\end{question}

\begin{question}[All-order inverse asymptotics]
Derive explicit functions $P_k(r)$ for a full expansion
$v_n\sim-S_n\ee^{-br_n}(1+P_1(r_n)/n+P_2(r_n)/n^2+\cdots)$,
with a precise meaning for the growing logarithmic coefficients.
The finite-core theorem supplies arbitrary algebraic tail savings, but this
paper does not perform the complete two-saddle expansion or claim the
polynomial form of every $P_k$ without proof.
\end{question}

\begin{question}[Which primitive data first enter each correction?]
The leading term depends only on $\lambda_1+w_2$. The next logarithmic core
jet is $2\lambda_1w_2+\tfrac32w_2^2-w_2\lambda_2-w_3$.
Determine the exact filtration of primitive data in successive inverse
asymptotic coefficients. Such a theorem would turn finite-core universality
into an accuracy-dependent information budget.
\end{question}

\begin{question}[Subquadratic feedback]
Extend the signed reduction to $\lambda_j=aj^p$ with $1<p<2$.
The quantity $jt$ at the inverse saddle then grows faster than logarithmically,
so a polynomial saving from a fixed crude majorant need not beat the core
cancellation. The quadratic proof identifies the obstruction rather than
solving it. A contour estimate that preserves more of the inverse core's
phase may be necessary.
\end{question}

\begin{question}[Regularly varying perturbations of the quadratic tail]
For $\lambda_j=aj^2L(j)$ with a controlled slowly varying function $L$,
determine the correct merger-defect estimate and inverse penalty.
The exact identity \eqref{eq:defect} is no longer available, but quantitative
convexity and regular variation may replace it. Mere pointwise
$\lambda_j/(aj^2)\to1$ is not enough to guarantee a multiplicative equivalent.
\end{question}

\begin{question}[Complex-direction Borel growth]
Determine an indicator-type description of $B(Re^{i\theta})$ on fixed and
shrinking angular sectors. The maximum-modulus theorem follows from eventual
one-sign coefficients, but does not identify the cancellation curves away
from the positive axis. Locate directions where ordinary Laplace growth
might hold and determine their boundaries.
\end{question}

\begin{question}[Acceleration at the Borel point at infinity]
Construct a summation or acceleration procedure adapted to
$\log(-B(R))\sim s\ee^{2s}/a$, $s\sim\sqrt{aR}$.
A useful result would identify a canonical analytic realization on a
specified sector and prove compatibility with formal inversion. The positive
ordinary Borel obstruction alone neither chooses such a procedure nor
proves uniqueness.
\end{question}

\begin{question}[A moving-scale asymptotic calculus]
Build an algebraic asymptotic calculus for sequences whose ratio is
$n/(\log n)^2$ times a constant, rather than a fixed factorial ratio.
The vanishing of every fixed positive-gamma-scale invariant in
\eqref{eq:gammaflat} suggests the need for additional moving-scale data.
Seek product and composition rules that recover the factor $\ee^{-br_n}$
without a separate saddle calculation in each example.
\end{question}

\begin{question}[Negative-real analytic realization]
For $t<0$, the Gaussian decay of $\ee^{aj^2t}$ makes the primitive tail
convergent for fixed $q$. Analyze the small solution of \eqref{eq:direct}
in complex neighborhoods of the negative axis, prove an asymptotic expansion
there, and compare it with any accelerated realization from the preceding
question. Domain control near $t=0$ must be included.
\end{question}

\begin{question}[Faster exact coefficient extraction]
The triangular integer algorithm is cubic in arithmetic operations.
Investigate fast truncated composition, divide-and-conquer recurrences,
and parameter-dependent polynomial multiplication to reduce its cost.
An effective implementation should report cancellation-sensitive working
precision separately from algebraic operation counts.
\end{question}

\begin{question}[A staged Lean formalization]
Begin with local finiteness, the frozen-parameter Lagrange identity,
the merger defect, and the ordered-composition majorant. Then formalize
analytic core estimates and the two saddle lemmas. These stages separate
finite algebra, inequalities, and complex analysis; none is represented
here as already checked by a proof assistant.
\end{question}

\section{Conclusion and proof boundary}
The signed quadratic inverse conjecture can be resolved without pretending
that its coefficients are positive. The method first sums the background
cancellation into an analytic inverse core, then pays for all remaining
large-action configurations with the exact quadratic merger defect.
A large auxiliary core makes that payment arbitrarily small on the relevant
polynomial scale. A separate one-tail saddle calculation then reduces the
leading answer to just two primitive actions.

The resulting equivalent is stable under arbitrary real finite perturbations,
with its leading exponential correction determined by $\lambda_1+w_2$.
It implies eventual negativity, a precise logarithmically damped factorial
growth scale, and an explicit entire Borel maximum-modulus equivalent.
The Borel obstruction occurs at infinity. None of these conclusions supplies
a general acceleration theorem, an optimal finite-$n$ certificate, a uniform
subquadratic theory, or a completed Lean development.

\appendix
\section{Core jets and the source of finite-order corrections}
\label{app:jets}
For $M\ge3$, expand $Q_M(t)=t-ct^2+q_3t^3+O(t^4)$.
Substitution into $H_M(Q_M(t),t)=t$ gives
\begin{align}
 c&=\lambda_1+w_2,\\
 q_3&=\frac{\lambda_1^2}{2}+3\lambda_1w_2+2w_2^2-w_2\lambda_2-w_3.
\end{align}
Hence
\begin{equation}
 \log\frac{Q_M(t)}t=-ct+
 \left(2\lambda_1w_2+\frac32w_2^2-w_2\lambda_2-w_3\right)t^2+O(t^3),
 \label{eq:logQ2}
\end{equation}
and
\begin{equation}
 A_M(t)=1-(\lambda_1+2w_2)t+O(t^2).
\end{equation}
In the pure model the quadratic logarithmic coefficient in
\eqref{eq:logQ2} is $\tfrac12-2a$.
At the inverse saddle its contribution has size
$j t^2=O((\log n)^3/n)$. This explains why the leading equivalent can
converge slowly while the two-core formula performs better at moderate
orders. These scale estimates are not a claimed complete first-correction
formula: the action-saddle displacement and its determinant also contribute.

\section{Proof-dependency audit}
\label{app:audit}
\begin{longtable}{>{\raggedright\arraybackslash}p{0.25\textwidth}>{\raggedright\arraybackslash}p{0.42\textwidth}>{\raggedright\arraybackslash}p{0.23\textwidth}}
\toprule Result & Essential inputs & Status\\\midrule\endhead
Formal inverse and marked response & Local finiteness; classical Lagrange inversion with frozen parameter & Proved here\\[5pt]
Bare saddle & Stirling; strict concavity; lattice Gaussian sum & Proved here\\[5pt]
Arbitrary-accuracy signed core reduction & Analytic core Cauchy estimate; exact merger defect; composition count; coarse localization & Proved here; constants not optimized\\[5pt]
Inverse equivalent & Core reduction plus uniform one-tail saddle & Proved here\\[5pt]
Minimal leading core & Same one-tail equivalent for every fixed $M\ge2$; explicit $M=1$ calculation & Proved here\\[5pt]
Finite signed perturbations & Fixed analytic core; unchanged eventual quadratic tail & Proved here; not uniform in the perturbation data\\[5pt]
Entire Borel equivalent & Coefficient theorem; eventual negativity; second real saddle & Proved here\\[5pt]
Forward/inverse ratio \eqref{eq:forward-comparison} & Repository forward theorem plus present inverse theorem & Explicit external dependency\\[5pt]
Finite exact checks & Integer recurrence; independent composition formulas; symbolic identities & Executed; not asymptotic proofs\\[5pt]
Lean verification, acceleration, nonreal Borel indicator & Additional work & Not supplied\\\bottomrule
\end{longtable}

The reference to a classical Lagrange formula is not a hidden appeal to a
positivity theorem. All sign-sensitive analytic estimates are displayed in
Sections~\ref{sec:reduction} and \ref{sec:one-tail}. The two instances of
asymptotic summation use positive comparison weights after the core
cancellation has already been extracted. No numerical observation is an
assumption in the proof chain.

\section{Notation guide}
\begin{longtable}{>{\raggedright\arraybackslash}p{0.24\textwidth}>{\raggedright\arraybackslash}p{0.68\textwidth}}
\toprule Symbol & Meaning\\\midrule\endhead
$\U,\V$ & Formal forward series and its compositional inverse.\\[4pt]
$v_n,V_n$ & Ordinary inverse coefficient and integer normalization $V_n=n!v_n$.\\[4pt]
$a,w_j,\lambda_j$ & Positive quadratic-tail constant; real primitive weights and slopes.\\[4pt]
$c,b,C_0$ & $\lambda_1+w_2$, $c/a$, and $1+|\lambda_1|+|w_2|$.\\[4pt]
$H_M,P_M,Q_M,A_M$ & Finite kernel; its inverse in the first variable; diagonal inverse; inverse derivative.\\[4pt]
$L_M,\ell_{n,M}$ & Formal one-tail inverse response and its ordinary coefficient; the inverse contains $-L_M$.\\[4pt]
$m,h_i,J,\Lambda$ & Number of tail actions; tail indices; their sum; sum of their slopes.\\[4pt]
$d_i,D,j,k,\Delta$ & Excess indices; total excess; merged action; coefficient complement; quadratic merger defect.\\[4pt]
$r_n,j_n,k_n,D_n,S_n$ & Bare saddle coordinate, action, complement, determinant factor, and coefficient equivalent.\\[4pt]
$B,R,s,n_R,\Sigma_R$ & Ordinary Borel transform; positive radius; explicit Borel saddle coordinate; index saddle; its width.\\\bottomrule
\end{longtable}
Here $D$ is total excess, $D_n$ a saddle determinant, and $\Delta$ the merger defect.

\begin{thebibliography}{9}
\bibitem{repo-core}
Vladimir Reshetnikov, ProveIt repository,
\emph{Finite-Core Universality and Sharp Large Order for Countable
Exponential-Feedback Transseries}, research manuscript, 29 September 2026.
Path: \path{Analysis/Transseries/docs/series-and-transseries/Finite_Core_Universality_Exponential_Feedback/article.tex}.
Inspected TeX blob: \path{c2b6faf623003111e9d9e6bbb162a67429f3ce42}.
\url{https://github.com/VladimirReshetnikov/ProveIt}.

\bibitem{gessel}
I. M. Gessel, \emph{Lagrange inversion}, Journal of Combinatorial Theory,
Series A \textbf{144} (2016), 212--249.
\url{https://doi.org/10.1016/j.jcta.2016.06.018};
\url{https://arxiv.org/abs/1609.05988}.

\bibitem{borinsky}
M. Borinsky, \emph{Generating asymptotics for factorially divergent sequences},
arXiv:1603.01236, version 4, 8 October 2018.
\url{https://arxiv.org/abs/1603.01236}.
The comparison uses the fixed positive-$\alpha$ gamma-scale definition and
its explicitly allowed all-zero asymptotic expansion.

\bibitem{sauzin}
D. Sauzin, \emph{Introduction to 1-summability and resurgence},
arXiv:1405.0356 (2014).
\url{https://arxiv.org/abs/1405.0356}.

\bibitem{edgar}
G. A. Edgar, \emph{Transseries for beginners}, arXiv:0801.4877 (2008).
\url{https://arxiv.org/abs/0801.4877}.

% ed. (2026-09-29): editorial bibliography entries for the repository articles cited in the editorial notes
\bibitem{ed:qef}
[Editorial addition] \emph{Quadratic Exponential Feedback after Reversion:
Cancellation, Finite-Core Universality, and Sharp Borel Growth}, research
article prepared for Vladimir Reshetnikov, ProveIt, filed in batch~49,
\path{Analysis/Transseries/docs/series-and-transseries/Quadratic_Exponential_Feedback_After_Reversion/article.tex}.
A second, independent claimed proof, not independently reviewed.

\bibitem{ed:nrs}
[Editorial addition] \emph{Negative-Ray Summation of Countable
Exponential-Feedback Transseries}, research article prepared for Vladimir
Reshetnikov, ProveIt, filed in batch~47 (in the inspected snapshot),
\path{Analysis/Transseries/docs/series-and-transseries/Negative_Ray_Summation_Exponential_Feedback/article.tex}.

\bibitem{ed:reg}
[Editorial addition] \emph{A sharp regularity classification for
countable exponential-feedback transseries}, research draft prepared for
Vladimir Reshetnikov, ProveIt, filed in batch~45 (in the inspected
snapshot),
\path{Analysis/Transseries/docs/series-and-transseries/Exponential_Feedback_Regularity_Classification/article.tex}.

\bibitem{ed:swt}
[Editorial addition] \emph{Sharp Weighted Type Under Formal Reversion},
research article prepared for Vladimir Reshetnikov, ProveIt, filed in
batch~48 (after the inspected snapshot),
\path{Analysis/Transseries/docs/series-and-transseries/Sharp_Weighted_Type_Formal_Reversion/article.tex}.
\end{thebibliography}
\end{document}
